% ===========================================================================
\section{Objective design}\label{sec:objectives}
% ===========================================================================

\subsection{Gradients with respect to student logits}

\begin{proposition}[Logit gradients of $f$-divergences]\label{prop:grad}
\Pv\Ck Let $q=\softmax(z)$, $u=p/q$, and $\psi_f(u):=f(u)-uf'(u)$. Then
$\nabla_zD_f(p\|q)=q\odot\big(\psi_f(u)-\E_q\psi_f(u)\big)$. In particular:
\begin{align*}
\nabla_z\KL(p\|q)&=q-p,\\
\nabla_z\KL(q\|p)&=q\odot\big(\log\tfrac qp-\KL(q\|p)\one\big),\\
\nabla_z\JSD_\beta(p\|q)&=(1-\beta)\,q\odot\big(\log\tfrac qm-\E_q\log\tfrac qm\,\one\big),\qquad m=\beta p+(1-\beta)q.
\end{align*}
\end{proposition}
\begin{proof}
$\partial_{q(y)}\sum_{y'}q(y')f(p(y')/q(y'))=f(u_y)-u_yf'(u_y)$. The softmax Jacobian is
$\partial q_j/\partial z_i=q_j(\delta_{ij}-q_i)$. For KL, $\psi=-u$. For reverse KL, written as $D_{-\log}$,
$\psi=1-\log u$. For $\JSD_\beta$ we have $f_\beta(u)-uf_\beta'(u)=-(1-\beta)\log(\beta u+1-\beta)=(1-\beta)\log(q/m)$.
\end{proof}

\subsection{Mode covering versus mode seeking}

\begin{proposition}[Support behaviour]\label{prop:support}
\Pv If $\KL(p\|q)<\infty$ then $\operatorname{supp}p\subseteq\operatorname{supp}q$ (zero-avoiding, mode covering). If
$\KL(q\|p)<\infty$ then $\operatorname{supp}q\subseteq\operatorname{supp}p$ (zero-forcing, mode seeking).
$\JSD_\beta\le h(\beta)$ is always finite and imposes neither constraint.
\end{proposition}

\begin{proposition}[Missing-mode gradients]\label{prop:missing}
\Pv Fix a token $i$ with $p_i>0$ and let $q_i\to0$, the other coordinates staying bounded away from $0$.
\begin{enumerate}[label=(\roman*),nosep]
\item Forward KL: $\partial_{z_i}\KL(p\|q)=q_i-p_i\to-p_i$. The restoring force equals the missing mass.
\item Reverse KL: $|\partial_{z_i}\KL(q\|p)|\le q_i\big(|\log\tfrac{q_i}{p_i}|+\KL(q\|p)\big)=O(q_i\log\tfrac1{q_i})\to0$.
A mode the student has dropped exerts \emph{no} pull.
\item $\JSD_\beta$: $\partial_{z_i}\JSD_\beta=-(1-\beta)\,q_i\big[\log(1-\beta+\beta p_i/q_i)+c\big]$ with
$c=\KL(q\|m)\in[0,\log\frac1{1-\beta}]$. The pull towards the missing mode is therefore at most
$\beta(1-\beta)p_i+(1-\beta)q_i\log\frac1{1-\beta}$, at most a quarter of the forward-KL pull when $q_i\ll p_i$.
As $q_i\to0$ it equals $(1-\beta)q_i\log\frac{\beta p_i}{q_i}(1+o(1))\to0$, which is reverse-like. A teacher mode is
effectively recovered only while the student's mass on it is not $\ll\beta p_i$.
\end{enumerate}
\end{proposition}
\begin{proof}
Read off \cref{prop:grad}. For (iii), $m_i/q_i=1-\beta+\beta p_i/q_i$, and
$0\le\E_q\log\frac qm\le\log\frac1{1-\beta}$ because $m\ge(1-\beta)q$. Use $q_i\log(1+\beta p_i/q_i)\le\beta p_i$ for the
upper bound, and $\log(1-\beta+\beta p_i/q_i)=\log\frac{\beta p_i}{q_i}+o(1)$ as $q_i\to0$.
\end{proof}

\begin{proposition}[M-projection versus I-projection]\label{prop:projections}
\Pv(a), \Kn(b).
(a)~For an exponential-family student $q_\eta\propto e^{\eta^\top\phi}$, $\eta\mapsto\KL(p\|q_\eta)$ is convex,
and its stationarity condition is moment matching, $\E_{q_\eta}\phi=\E_p\phi$. Stationarity of
$\KL(q_\eta\|p)$ is $\Cov_{q_\eta}(\phi,\log\frac{q_\eta}p)=0$, which is generally non-convex.
(b)~Let $p=w\,\cN(-a,1)+(1-w)\,\cN(a,1)$ with $w<\frac12$, and let $q=\cN(\mu,s^2)$.
The forward minimiser is $\mu=(1-2w)a$, $s^2=1+4w(1-w)a^2$. At it, $\KL(q\|p)=\Theta(a^2)$ and $\TV(p,q)\to1$
as $a\to\infty$.
The reverse KL has local minimisers near $\mu=\pm a$ with $s\to1$, of value $\to\log\frac1w$ and
$\log\frac1{1-w}$ respectively; the global one is the heavier mode. At it, $\KL(p\|q)=\Theta(a^2)$ but
$\TV(p,q)\to w$.
\end{proposition}
\begin{proof}
(a) $\nabla_\eta\KL(p\|q_\eta)=\E_{q_\eta}\phi-\E_p\phi$, and the Hessian is $\Cov_{q_\eta}\phi\succeq0$. For the
reverse direction, differentiate $\E_{q_\eta}[\log q_\eta-\log p]$.
(b) The forward case is moment matching. The mixture has mean $(1-2w)a$ and variance
$1+a^2-(1-2w)^2a^2$. The reverse values follow because $p\approx(1-w)\cN(a,1)$ on the bulk of
$\cN(a,1)$. Classical discussions are \citet{minka2005divergence} and \citet{bishop2006pattern}.
\end{proof}
\noindent
Each direction's minimiser is catastrophic in the other direction. In this example, the mode-seeking
solution is \emph{behaviourally} better (TV $\approx w$ versus $\approx1$): it produces valid samples and
drops a mode, whereas the covering solution produces mostly invalid samples. This is the formal content
of ``mode covering versus mode seeking''. Which is preferable depends on the deployment criterion
(\cref{sec:right}). Under misspecification, neither is TV-optimal (\cref{prop:agnostic}).

\subsection{Temperature}

Let $p^\tau=\softmax(v/\tau)$, $q^\tau=\softmax(z/\tau)$, and $\cL_\tau(z):=\tau^2\KL(p^\tau\|q^\tau)$, the loss of
\citet{hinton2015distilling}.
\begin{proposition}[Temperature and the $\tau^2$ factor]\label{prop:temperature}
\Pv\Ck
\begin{enumerate}[label=(\roman*),nosep]
\item $\nabla_z\cL_\tau=\tau(q^\tau-p^\tau)$. Without the factor $\tau^2$ the gradient is $\frac1\tau(q^\tau-p^\tau)=O(\tau^{-2})$
as $\tau\to\infty$. The factor $\tau^2$ keeps gradient magnitudes $O(1)$, so the relative weighting against a
hard-label term is stable across $\tau$.
\item (High temperature: logit matching.)
$\tau(q^\tau_i-p^\tau_i)=\frac1V(\tilde z_i-\tilde v_i)+\frac1{2V\tau}\big[(\tilde z_i^2-\tilde v_i^2)-\frac1V(\|\tilde z\|^2-\|\tilde v\|^2)\big]+O(\tau^{-2})$,
so $\cL_\tau(z)=\frac1{2V}\|\tilde z-\tilde v\|^2+O(1/\tau)$. This is centred-logit regression, as in
\citet{ba2014do}.
\item (Low temperature: argmax matching.) If $v$ has a unique maximiser $i^\star$, then
$\tau\KL(p^\tau\|q^\tau)\to\max_jz_j-z_{i^\star}$, a perceptron/hinge loss on the teacher's argmax, and $\cL_\tau\to0$.
\item (Geometry, not solution.) The global minimisers $\{v+c\one\}$ do not depend on $\tau$, and
$\nabla^2\cL_\tau|_{z=v}=\diag(p^\tau)-p^\tau p^{\tau\top}\to\frac1VC$ as $\tau\to\infty$. Temperature changes \emph{which
errors are cheap}, and hence the misspecified solution and the conditioning, but not the realisable
solution.
\end{enumerate}
\end{proposition}
\begin{proof}
(i) Chain rule: $\partial_z\KL(p^\tau\|\softmax(z/\tau))=\frac1\tau(q^\tau-p^\tau)$.
(ii) Expand $e^{\tilde z_i/\tau}=1+\tilde z_i/\tau+\tilde z_i^2/2\tau^2+O(\tau^{-3})$, with
$\sum_je^{\tilde z_j/\tau}=V+\|\tilde z\|^2/2\tau^2+O(\tau^{-3})$, and similarly for $v$.
(iii) $p^\tau\to\delta_{i^\star}$, and $-\tau\log q^\tau_{i^\star}=\tau\,\mathrm{lse}(z/\tau)-z_{i^\star}\to\max_jz_j-z_{i^\star}$;
the entropy terms of $p^\tau$ vanish.
(iv) The Hessian of $\tau^2\mathrm{lse}(z/\tau)$ is $\diag(q^\tau)-q^\tau q^{\tau\top}$.
\end{proof}

\subsection{Dark knowledge as Fisher information}

\begin{proposition}[Low-probability tokens are invisible at $\tau=1$]\label{prop:invisible}
\Pv For $\delta\in\R^V$ supported on $A\subseteq\cV$,
$\delta^\top(\diag p-pp^\top)\delta\le\max_{i\in A}p_i\,\|\delta\|^2\le p(A)\|\delta\|^2$. Near the optimum,
perturbing the logits of tokens of total teacher mass $10^{-6}$ by $\delta$ costs at most $\frac12\cdot10^{-6}\|\delta\|^2$.
Their relative logits, which is the ``dark knowledge'', are essentially unconstrained at $\tau=1$. At
temperature $\tau$, $p$ is replaced by the flatter $p^\tau$.
\end{proposition}

\begin{proposition}[What soft labels add to hard labels]\label{prop:soft-info}
\Pv
\begin{enumerate}[label=(\alph*),nosep]
\item At $\theta^\circ$, one hard label at $s$ carries Fisher information $\mathcal I(s)=J^\top(\diag p-pp^\top)J$,
and $m$ i.i.d.\ hard labels carry $m\mathcal I(s)$. The soft label is the $m\to\infty$ limit. It determines the
centred teacher logits $\tilde v(s)$ exactly, because $p\mapsto C\log p$ is a bijection from the open simplex
onto $\operatorname{range}C$.
\item The per-state gradient-variance reduction from using the soft label is $\tr\mathcal I(s)$
(\cref{prop:rb}). It is zero iff the teacher is deterministic at $s$ (for $J$ of full rank on
$\operatorname{range}C$). In the realisable linear case it upgrades a $k/2n$ error to exact recovery
(\cref{prop:exact}).
\item For any $\tau>0$, $p^\tau$ is a bijective function of $p$. Temperature adds \emph{no information}; it
reweights existing information in the loss (\cref{prop:temperature}(iv), \cref{prop:invisible}). It
matters only when the channel is lossy, for example with top-$k$ logits, finite precision, or finite
optimisation budgets.
\item Under \cref{as:linear} for the teacher, $\tilde v(s)=CW_\TT h_\TT(s)$. The soft label therefore reveals
the teacher's final hidden state up to $\ker(CW_\TT)$: dark knowledge is a linear image of the
teacher's representation. This links to \cref{sec:repr}.
\end{enumerate}
\end{proposition}
\noindent
So soft labels carry information beyond hard labels \emph{exactly when the teacher is non-degenerate}.
The amount is unbounded in the noiseless realisable case and equals the Rao--Blackwell gap otherwise.
Temperature is a preconditioner, not an information source. This sharpens the ``dark knowledge'' account
of \citet{hinton2015distilling} and the gradient decompositions of \citet{furlanello2018born} and
\citet{tang2020understanding}. Relations to label smoothing are discussed by \citet{muller2019does}.

\subsection{Sequence-level objectives and unbiased gradient estimators}

\begin{proposition}[Sequence-level $f$-divergence gradients]\label{prop:seq-grad}
\Pv\Ck For sequence laws $P$ and $Q_\theta$,
$\nabla_\theta D_f(P\|Q_\theta)=\E_{y\sim Q_\theta}\big[\nabla_\theta\log Q_\theta(y)\,\psi_f\big(P(y)/Q_\theta(y)\big)\big]$.
This gives
$\nabla\KL(Q_\theta\|P)=\E_{Q_\theta}[\nabla\log Q_\theta\cdot\log\frac{Q_\theta}P]$,
$\nabla\JSD_\beta(P\|Q_\theta)=(1-\beta)\E_{Q_\theta}[\nabla\log Q_\theta\cdot\log\frac{Q_\theta}{M_\beta}]$ with $M_\beta$
held fixed inside the logarithm, and $\nabla\KL(P\|Q_\theta)=-\E_P[\nabla\log Q_\theta]$ (off-policy).
\end{proposition}
\begin{proof}
$\nabla D_f=\sum_y\nabla Q_\theta(y)\,\partial_{Q(y)}D_f$ with $\partial_{Q(y)}D_f=\psi_f(u_y)$ as in \cref{prop:grad}.
Constants integrate to zero against $\nabla Q_\theta$. The divergence-minimisation view of imitation
learning is developed by \citet{ghasemipour2019divergence} and \citet{ke2020imitation}, and for sequence-level
distillation by \citet{wen2023fdivergence}.
\end{proof}

\begin{proposition}[Rao--Blackwellised unbiased estimator for sequence-level reverse KL]\label{prop:rb-reverse}
\Pv\Ck Let $y\sim Q^x_\theta$, $g_t(s):=\KL(q_\theta(\cdot\mid s)\|p(\cdot\mid s))$, and let $b_t$ be any function of $s_t$
(a baseline). Then
\[
\hat g=\sum_{t=1}^T\Big[\nabla_\theta g_t(s_t)\big|_{s_t\ \text{fixed}}+\nabla_\theta\log q_\theta(y_t\mid s_t)\Big(\sum_{k>t}g_k(s_k)-b_t(s_t)\Big)\Big]
\]
satisfies $\E\hat g=\nabla_\theta\KL(Q^x_\theta\|\Pseq)$. The ``stop-gradient'' estimator that keeps only the first
term is \emph{biased}. Its bias is $-\sum_t\E[\nabla\log q_\theta(y_t\mid s_t)\sum_{k>t}g_k(s_k)]$, the effect of the
student's actions on its own future prefixes. Under \cref{as:real} both estimators vanish in expectation
at $\theta^\circ$. In general the stop-gradient fixed points are those of the DAgger-type map
$\theta\mapsto\argmin_{\theta'}\sum_t\E_{\dstud t}\KL(q_{\theta'}\|p)$, not stationary points of the sequence-level KL.
\end{proposition}
\begin{proof}
$F(\theta)=\sum_t\E_{s\sim\dstud t}g_t(\theta;s)$ by \cref{thm:chain}. Differentiating gives a direct term and a
score term through $\dstud t$:
$\sum_t\E[(\sum_{k<t}\nabla\log q(y_k\mid s_k))g_t(s_t)]=\sum_k\E[\nabla\log q(y_k\mid s_k)\sum_{t>k}g_t(s_t)]$.
Replacing the sampled per-token log-ratios by $g_k$ is exact by the tower property, since
$\E[\log\frac{q}{p}(y_k\mid s_k)\mid s_k]=g_k(s_k)$, and $s_k$ contains $y_t$ for $k>t$. Baselines are unbiased
because $\E[\nabla\log q(y_t\mid s_t)\mid s_t]=0$. The analytic per-token term is itself the
Rao--Blackwellisation of $\nabla\log q(y_t\mid s_t)\log\frac q p(y_t\mid s_t)$.
\citet{gu2024minillm} use a policy gradient for this objective (with further variance-reduction
devices). \citet{agarwal2024gkd} do not differentiate through the student's sampling; we believe this is
their stated choice, but we have not re-verified the exact wording.
\end{proof}

\begin{proposition}[Sequence-level $\JSD_\beta$ decomposes with prefix-adaptive mixtures]\label{prop:pajsd}
\Pv The mixture $M_\beta=\beta P+(1-\beta)Q_\theta$ is autoregressive, with conditionals
\[
m_t(\cdot\mid s)=w_t(s)\,p_t(\cdot\mid s)+(1-w_t(s))\,q_t(\cdot\mid s),\qquad
w_t(s)=\sigma\Big(\log\tfrac\beta{1-\beta}+\log P(s)-\log Q_\theta(s)\Big),
\]
where $P(s)=\prod_{k<t}p(y_k\mid s_k)$ and $Q_\theta(s)$ are prefix probabilities. Consequently
\[
\JSD_\beta(\Pseq\|\Qseq)=\beta\sum_{t=1}^T\E_{s\sim\Pseq}\KL\big(p_t(s)\,\|\,m_t(s)\big)+(1-\beta)\sum_{t=1}^T\E_{s\sim\Qseq}\KL\big(q_t(s)\,\|\,m_t(s)\big).
\]
The token-level $\JSD_\beta$ of \citet{agarwal2024gkd} is the case $w_t\equiv\beta$. It coincides with the
sequence-level divergence on prefixes with $P(s)=Q_\theta(s)$, so it is its \emph{near-lossless linearisation}.
Because $\log P(s_t)-\log Q_\theta(s_t)$ is a sum of $t-1$ per-token log-ratios, the two diverge along long
prefixes unless the student is already nearly lossless there.
\end{proposition}
\begin{proof}
$M_\beta(y_{\le t})=\beta P(y_{\le t})+(1-\beta)Q_\theta(y_{\le t})$. Dividing by $M_\beta(y_{<t})$ gives $m_t$. Now
apply the chain rule (\cref{thm:chain}), which holds for any pair of autoregressive laws, to $\KL(P\|M_\beta)$
and to $\KL(Q_\theta\|M_\beta)$.
\end{proof}

\begin{proposition}[Unbiased per-token estimator for sequence-level $\JSD_\beta$]\label{prop:jsd-est}
\Pv Let $y\sim\Qseq$, let $\rho_t(a;s):=\log\frac{q_t(a\mid s)}{m_t(a\mid s)}$ be evaluated with \emph{stop-gradient}
on $m_t$ and $w_t$, and let $b_t$ be any function of $s_t$. Then
\begin{multline*}
\hat g_{\JSD}=(1-\beta)\sum_{t=1}^T\Big[\sum_{a\in\cV}\nabla_\theta q_\theta(a\mid s_t)\,\rho_t(a;s_t)\\
+\nabla_\theta\log q_\theta(y_t\mid s_t)\Big(\sum_{k>t}\KL\big(q_k(s_k)\|m_k(s_k)\big)-b_t(s_t)\Big)\Big]
\end{multline*}
is unbiased for $\nabla_\theta\JSD_\beta(\Pseq\|\Qseq)$. It uses student samples only. The teacher enters through $m_t$.
\end{proposition}
\begin{proof}
By \cref{prop:seq-grad} and \cref{prop:pajsd},
$\nabla\JSD_\beta=(1-\beta)\E_Q[\sum_t\nabla\log q(y_t\mid s_t)\sum_k\rho_k(y_k;s_k)]$. Terms with $k<t$ vanish, because
$\rho_k$ is $s_t$-measurable. The $k=t$ term is replaced by its conditional expectation over $y_t$, and
$k>t$ terms by $\E[\rho_k\mid s_k]=\KL(q_k\|m_k)$. Baselines are unbiased as before.
\end{proof}

\subsection{Bounded divergences are agnostically TV-competitive; KL is not}

\begin{proposition}[Agnostic TV-competitiveness]\label{prop:agnostic}
\Pv
\begin{enumerate}[label=(\alph*),nosep]
\item For any student class, every $\theta_\beta\in\argmin_\theta\E_x\JSD_\beta(\Pseq\|\Qseq)$ satisfies
\[
\E_x\TV(\Pseq,Q^x_{\theta_\beta})\le\sqrt{\frac{h(\beta)}{2\beta(1-\beta)}\;\inf_\theta\E_x\TV(\Pseq,\Qseq)} .
\]
The constant is minimised at $\beta=\frac12$, giving $\sqrt{2\ln2\cdot\mathrm{TV}^\star}\approx1.18\sqrt{\mathrm{TV}^\star}$.
\item No bound of the form $\TV(\theta_{\KL})\le\phi(\inf_\theta\TV)$ with $\phi(0^+)=0$ holds for either KL
direction. Let $P=(1-\eta)\delta_a+\eta\delta_b$.
\emph{Forward:} for the class $\{\delta_a,\mathrm{Unif}(\cV)\}$, the KL minimiser is the uniform law, with
$\TV\ge1-\eta-\frac1V$, while $\inf\TV=\eta$.
\emph{Reverse:} for the class $\{(1-\eta')\delta_a+\eta'\delta_c,\ \delta_b\}$ with $c\notin\operatorname{supp}P$, the
reverse-KL minimiser is $\delta_b$ (value $\log\frac1\eta$), with $\TV=1-\eta$, while $\inf\TV\le\eta+\eta'$.
\end{enumerate}
\end{proposition}
\begin{proof}
(a) Let $\theta_\TV$ attain (or approach) $\inf\E\TV$. Using \cref{prop:jsd}(i) twice and Jensen,
\[
\E\TV(\theta_\beta)\le\sqrt{\frac{\E\JSD_\beta(\theta_\beta)}{2\beta(1-\beta)}}\le\sqrt{\frac{\E\JSD_\beta(\theta_\TV)}{2\beta(1-\beta)}}\le\sqrt{\frac{h(\beta)\,\E\TV(\theta_\TV)}{2\beta(1-\beta)}}.
\]
The function $h(\beta)/\beta(1-\beta)$ is symmetric about $\frac12$ and is minimised there (direct check).
(b) Direct computation. $\KL(P\|\delta_a)=\infty$, and $\KL(\cdot\|P)=\infty$ for any law charging $c$.
\end{proof}

\noindent
\Cref{prop:agnostic} is the formal reason to prefer a bounded divergence when the student is
capacity-limited. KL-based objectives can place their unavoidable error where TV, and hence every bounded
downstream task, is maximally hurt.

\subsection{Proposed objective}\label{sec:proposal}

\begin{definition}[Prefix-adaptive sequence-level JSD distillation, PA-JSD]\label{def:pajsd}
\[
\cJ_\beta(\theta):=\E_{x\sim\cD}\JSD_\beta\big(\Pseq\,\|\,\Qseq\big),\qquad\beta=\tfrac12\text{ by default},
\]
computed through the exact decomposition of \cref{prop:pajsd}. The first sum is estimated on
teacher-sampled prefixes with the teacher's soft labels. The second is estimated on student roll-outs.
Gradients use \cref{prop:jsd-est}, optionally plus the pathwise gradient of the teacher-prefix term.
Optionally add $\gamma\,\cR_{\mathrm{F}}$, the Fisher-aligned representation term of \cref{def:fisher-rep}.
\end{definition}

\begin{theorem}[Properties of PA-JSD]\label{thm:pajsd}
\Pv (except where marked)
\begin{enumerate}[label=(\roman*),nosep]
\item \emph{Exactness (no exposure bias).} $\cJ_\beta$ is a sequence-level divergence, and its per-token form
is an identity (\cref{prop:pajsd}). With $\cE=\E\JSD_\beta$, \cref{thm:master} holds with
$\cE_{\mathrm{shift}}\equiv0$.
\item \emph{Behavioural control.} $\E_\cD\TV(P,Q_\theta)\le\sqrt{\cJ_\beta(\theta)/(2\beta(1-\beta))}$. Under
input shift, $\E_\mu\JSD_\beta\le\cJ_\beta+h(\beta)\TV(\mu,\cD)$; no density-ratio assumption is needed.
\item \emph{Agnostic competitiveness.} The population minimiser is TV-competitive (\cref{prop:agnostic}).
\item \emph{Fisher consistency.} $\cJ_\beta(\theta)=0$ iff $\Qseq=\Pseq$ for $\cD$-a.e.\ $x$. Under \cref{as:real},
the minimisers are exactly the lossless students.
\item \emph{No loss in the near-lossless regime.} If $a\le\Pseq/\Qseq\le b$, then
$\cJ_\beta=\frac{\beta(1-\beta)}2\E\chi^2(1+O(b-a))=\beta(1-\beta)\,\E\KL(P\|Q_\theta)(1+O(b-a))$.
\item \emph{Unbiased, Rao--Blackwellised gradients} (\cref{prop:jsd-est}).
\item \emph{Horizon-free estimation} \Pa{\cref{as:finite}}. Per-prompt values lie in $[0,h(\beta)]$.
If the per-prompt divergences are evaluated exactly, then with probability $\ge1-\delta$ the empirical
objective deviates uniformly by at most $h(\beta)\sqrt{(R\ln2+\log(2/\delta))/(2n)}$, with no factor $T$.
Compare the factor $TB$ in \cref{thm:unif}.
\end{enumerate}
\end{theorem}
\begin{proof}
(i) \cref{prop:pajsd}. (ii) \cref{prop:jsd}(i) with Jensen, and \cref{prop:covshift} with $G=h(\beta)$.
(iii) \cref{prop:agnostic}. (iv) $\JSD_\beta\ge2\beta(1-\beta)\TV^2$. (v) \cref{prop:local}(ii) applied to
sequence laws, and $f_\beta''(1)=\beta(1-\beta)$. (vi) \cref{prop:jsd-est}. (vii) Hoeffding for $[0,h(\beta)]$-valued
losses and a union bound over $2^R$ students.
\end{proof}

\begin{remark}[Why this targets the dominant term]
By \cref{heur:dominant}, the dominant term under capacity limits is approximation error as seen by the
deployment criterion. PA-JSD attacks all three ways that term is inflated. Exactness removes its
amplification by prefix shift. Agnostic competitiveness prevents its misallocation. Boundedness makes
the estimation and covariate-shift terms horizon-free and additive. In the near-lossless regime it
reduces to KL up to the constant $\beta(1-\beta)$, so nothing is lost when losslessness is reachable. The
price is on-policy sampling and the variance of a sequence-level score function. The saturation
$w_t\in\{0,1\}$ on prefixes that are already distinguishable limits the effective horizon of that
variance, but does not remove it.
\end{remark}

\begin{proposition}[Likelihood-critical alternative: skew-Jeffreys]\label{prop:sj}
\Pv Let $\cJ^{\mathrm{SJ}}_\lambda(\theta):=(1-\lambda)\sum_t\E_{\dteach t}\KL(p_t\|q_t)+\lambda\sum_t\E_{\dstud t}\KL(q_t\|p_t)$. It is
estimated by off-policy soft-label forward KL plus on-policy reverse KL (\cref{prop:rb-reverse}). Then
$\cJ^{\mathrm{SJ}}_\lambda=\E[(1-\lambda)\KL(P\|Q_\theta)+\lambda\KL(Q_\theta\|P)]$ exactly, and
$\E\TV\le\sqrt{\cJ^{\mathrm{SJ}}_\lambda/2}$. The log-loss regret is at most $\cJ^{\mathrm{SJ}}_\lambda/(1-\lambda)$.
Both event bounds of \cref{prop:events} hold with $\KL$ replaced by $\cJ^{\mathrm{SJ}}_\lambda/(1-\lambda)$
(coverage) and $\cJ^{\mathrm{SJ}}_\lambda/\lambda$ (precision), per input. It is not agnostically
TV-competitive (\cref{prop:agnostic}(b)).
\end{proposition}
\noindent
\textbf{Recommendation.} If two-sided losslessness is within reach, the two objectives agree to second
order (\cref{prop:local}), and skew-Jeffreys additionally \emph{certifies} likelihood-level guarantees.
If the student is capacity-limited, use PA-JSD to allocate its error where the deployment criterion is
least hurt. A natural schedule is PA-JSD to convergence followed by a skew-Jeffreys fine-tune, monitoring
both sequence-level quantities (\cref{sec:experiments}).

